\documentclass[11pt]{article}
\usepackage[margin=1in]{geometry}
\usepackage{amsmath,amssymb,amsthm,mathtools}
\usepackage{graphicx}
\usepackage{microtype}
\usepackage[hidelinks]{hyperref}

\newtheorem{theorem}{Theorem}
\newtheorem{proposition}{Proposition}
\newtheorem{remark}{Remark}
\newcommand{\R}{\mathbb R}
\newcommand{\K}{\mathcal K}
\newcommand{\Lop}{\mathcal L}

\title{A Minimal Three-Dimensional Counterexample for Eventual Cone Invariance\\
\large A candidate full answer to the question after Example 3.1 of arXiv:2303.07809}
\author{Autonomous counterexample candidate for human review}
\date{11 August 2026}

\begin{document}
\maketitle

\begin{center}
\fbox{\parbox{0.91\textwidth}{\textbf{Status.}
Candidate full counterexample and minimal-dimension resolution, likely valid.
An explicit nilpotent $3\times3$ matrix is individually but not uniformly
eventually nonnegative on the $2\times2$ positive-semidefinite cone.  A bonus
restriction argument also improves a secondary dimension bound from $16$ to
$6$.}}
\end{center}

\section{Source question}

Glück and Hölz \cite{GH} construct a four-dimensional ordered space and a
matrix semigroup that is individually but not uniformly eventually
nonnegative.  Immediately after the construction they ask whether a
counterexample exists in dimension three.

\begin{figure}[ht]
 \centering
 \includegraphics[width=\textwidth]{figures/open_question_crop.png}
 \caption{The dimension-three question after Example 3.1, source PDF page 7.}
\end{figure}

Recall that individual eventual nonnegativity means that for each $x$ in the
cone there is a time $t_0(x)$ such that $e^{tA}x$ stays in the cone for every
$t\geq t_0(x)$.  Uniform eventual nonnegativity requires one common $t_0$ for
all cone vectors.

\section{The three-dimensional counterexample}

Let
\begin{equation}\label{eq:cone}
 \K=\{(x_1,x_2,x_3)\in\R^3:
 x_1\geq0,\ x_3\geq0,\ x_2^2\leq x_1x_3\}.
\end{equation}
Under the identification
\[
 (x_1,x_2,x_3)\longleftrightarrow
 \begin{pmatrix}x_1&x_2\\x_2&x_3\end{pmatrix},
\]
this is the positive-semidefinite cone.  It is therefore a closed, pointed,
convex cone with nonempty interior.

Set
\begin{equation}\label{eq:A}
 A=\begin{pmatrix}
 0&4&0\\0&0&1\\0&0&0
 \end{pmatrix}.
\end{equation}
Since $A^3=0$,
\begin{equation}\label{eq:exp}
 e^{tA}=I+tA+\frac{t^2}{2}A^2
 =\begin{pmatrix}
 1&4t&2t^2\\0&1&t\\0&0&1
 \end{pmatrix}.
\end{equation}

\begin{theorem}\label{thm:main}
The semigroup $(e^{tA})_{t\geq0}$ is individually eventually nonnegative with
respect to $\K$, but is not uniformly eventually nonnegative.  Indeed,
$e^{tA}\K\not\subseteq\K$ for every $t>0$.
\end{theorem}

\begin{proof}
Fix $x=(x_1,x_2,x_3)\in\K$ and write $y=e^{tA}x$.  Formula
\eqref{eq:exp} gives
\[
 y_1=x_1+4tx_2+2t^2x_3,\qquad
 y_2=x_2+tx_3,\qquad y_3=x_3.
\]
A direct expansion yields the determinant identity
\begin{equation}\label{eq:detidentity}
 y_1y_3-y_2^2
 =(x_1x_3-x_2^2)+2t x_2x_3+t^2x_3^2.
\end{equation}

If $x_3=0$, the cone inequality forces $x_2=0$, so
$e^{tA}x=x$ for every $t\geq0$.  Suppose $x_3>0$.  Since the first term on
the right of \eqref{eq:detidentity} is nonnegative, the whole expression is
nonnegative whenever
\[
 t\geq t_0(x):=\max\left\{0,-\frac{2x_2}{x_3}\right\}.
\]
For these times $y_3=x_3>0$, and the determinant inequality also implies
$y_1\geq y_2^2/y_3\geq0$.  Thus $e^{tA}x\in\K$ for every
$t\geq t_0(x)$.  This proves individual eventual nonnegativity.

Now fix an arbitrary $t>0$ and take
\begin{equation}\label{eq:witness}
 x(t)=(t^2,-t,1).
\end{equation}
It lies on the boundary of $\K$, because $(-t)^2=t^2\cdot1$.  But
\[
 e^{tA}x(t)=(-t^2,0,1)\notin\K.
\]
Hence $e^{tA}$ is not cone-preserving at any positive time, which rules out
uniform eventual nonnegativity.
\end{proof}

\begin{proposition}\label{prop:minimal}
Dimension three is minimal for an individual-versus-uniform separation.
\end{proposition}

\begin{proof}
A closed pointed cone with nonempty interior in a two-dimensional real vector
space is generated by its two boundary rays, so it is polyhedral.  If $u,v$
generate the cone and a semigroup is individually eventually nonnegative,
choose eventual times $t_u,t_v$ for those two vectors.  For
$t\geq\max\{t_u,t_v\}$ and every $\alpha,\beta\geq0$,
\[
 e^{tA}(\alpha u+\beta v)
 =\alpha e^{tA}u+\beta e^{tA}v
\]
is again in the cone.  Thus the eventual time is uniform.  The same is
immediate in dimension one.  Theorem~\ref{thm:main} therefore occurs in the
smallest possible dimension.
\end{proof}

\section{Bonus: weak versus individual in dimension six}

The source derives a weakly but not individually eventually nonnegative
semigroup on the $16$-dimensional operator space over its four-dimensional
example.  Theorem~\ref{thm:main} first reduces that construction to dimension
$9$.  It can be restricted further.

Let $X=\R^3$ with cone $\K$, and let $\mathcal C=\Lop(X)_+$ denote the cone of
$\K$-preserving operators.  Choose
\[
 u=(1,0,1)\in\operatorname{int}\K,\qquad
 \varphi(x)=x_1+x_3\in\operatorname{int}\K',\qquad
 U=u\otimes\varphi.
\]
Thus $U(\K\setminus\{0\})\subset\operatorname{int}\K$, so $U$ is an
interior point of $\mathcal C$.  Define
\begin{equation}\label{eq:Y}
 Y=\operatorname{span}\{I,A,A^2,U,AU,A^2U\},
 \qquad Y_+=Y\cap\mathcal C.
\end{equation}
The six displayed matrices are linearly independent: in matrix-unit notation,
\[
\begin{array}{lll}
 A=4E_{12}+E_{23},&A^2=4E_{13},&
 U=E_{11}+E_{13}+E_{31}+E_{33},\\
 AU=E_{21}+E_{23},&A^2U=4E_{11}+4E_{13},&I=E_{11}+E_{22}+E_{33}.
\end{array}
\]
Their unique matrix-unit entries successively force all coefficients in a
linear relation to vanish.  Hence $\dim Y=6$.  Moreover $Y_+$ is a closed cone
with nonempty relative interior because it contains the ambient interior point
$U$.

Let $B:Y\to Y$ be left multiplication by $A$.  The space is $B$-invariant,
since the two triples in \eqref{eq:Y} are nilpotent $A$-chains, and
\begin{equation}\label{eq:left}
 e^{tB}M=e^{tA}M.
\end{equation}

\begin{proposition}\label{prop:six}
The six-dimensional ordered-space semigroup $(e^{tB})_{t\geq0}$ is weakly but
not individually eventually nonnegative.
\end{proposition}

\begin{proof}
Because $Y$ contains an interior point of $\mathcal C$, the finite-dimensional
positive-extension theorem extends every positive functional on $(Y,Y_+)$ to
a positive functional on $(\Lop(X),\mathcal C)$.  The latter dual cone is the
conical hull of the evaluation functionals
\[
 M\longmapsto\langle x',Mx\rangle,\qquad x\in\K,\quad x'\in\K';
\]
this is Proposition 2.3 of \cite{GH}.

Fix $M\in Y_+$.  For each such evaluation, $Mx\in\K$, so
Theorem~\ref{thm:main} gives
\[
 \langle x',e^{tB}Mx\rangle
 =\langle x',e^{tA}Mx\rangle\geq0
\]
for all sufficiently large $t$.  A positive functional is a finite conical
combination of evaluations, so one common maximum of their individual times
proves weak eventual nonnegativity on $Y$.

On the other hand $I\in Y_+$, while
$e^{tB}I=e^{tA}\notin\mathcal C$ for every $t>0$ by
\eqref{eq:witness}.  Thus the positive vector $I$ never has an eventually
nonnegative orbit in $Y$, and the semigroup is not individually eventually
nonnegative.
\end{proof}

No claim is made that six is the smallest dimension for this second
separation.

\section{Verification and novelty status}

The main construction uses only a nilpotent exponential and one quadratic
cone inequality.  A bounded search of the run indexes, exact phrases from the
source, the journal citation page, and later web/arXiv records found no prior
dimension-three construction.  Priority is not asserted; both results should
receive specialist review before dissemination.

\begin{thebibliography}{9}
\bibitem{GH}
J.~Glück and J.~Hölz,
\emph{Eventual cone invariance revisited},
Linear Algebra Appl. 675 (2023), 274--293;
arXiv:2303.07809.
\end{thebibliography}

\end{document}
